\documentclass[12pt]{article}
\usepackage[russian]{babel}
\usepackage[utf8]{inputenc}
\usepackage{amsmath, amssymb, amsthm}
\usepackage{geometry}
\usepackage{graphicx}
\usepackage{hyperref}
\usepackage{cite}

\geometry{margin=1in}

\title{Электрон как прецессирующий магнитный диполь: \\ Классическая модель и связь с опытом Штерна-Герлаха}
\author{Алексей Дроздов}
\date{27.02.2026}

\begin{document}

\maketitle

\begin{abstract}
В данной статье представлена классическая модель электрона как прецессирующего магнитного диполя со внутренней структурой. Модель решает фундаментальную проблему спина, не требуя сверхсветовых скоростей вращения. Мы выводим уравнения движения прецессирующего диполя в неоднородных магнитных полях и демонстрируем, как эта модель объясняет дискретные результаты, наблюдаемые в опыте Штерна-Герлаха. Анализ выявляет дополнительные стационарные состояния помимо традиционных двух, и устанавливает соответствие между классическими параметрами и квантово-механическим формализмом (уравнение Дирака, матрицы Паули). Модель предоставляет механистическую интерпретацию спина, избегая концептуальных трудностей квантовой механики точечных частиц.
\end{abstract}

\section{Введение}

Опыт Штерна-Герлаха (1922) продемонстрировал, что атомные пучки расщепляются на дискретные компоненты при прохождении через неоднородные магнитные поля. Этот результат фундаментально противоречил классической электродинамике, которая предсказывала непрерывное распределение отклонений. Квантовая механика разрешила этот парадокс, постулируя собственный угловой момент (спин) с квантованными проекциями.

Однако квантово-механическое описание спина остаётся чисто формальным — оно постулирует, а не объясняет. Уравнение Дирака выводит спин естественным образом из релятивистской инвариантности, но не даёт механической картины того, что такое спин на самом деле.

В данной статье предлагается классическая модель, где электрон рассматривается как прецессирующий магнитный диполь со внутренней структурой. Этот подход:
\begin{itemize}
\item Избегает сверхсветовых скоростей вращения, которые преследуют классические модели вращающейся сферы
\item Объясняет два дискретных состояния, наблюдаемых в опытах Штерна-Герлаха
\item Предсказывает дополнительные стационарные состояния при определённых условиях
\item Устанавливает соответствие с квантово-механическим формализмом
\end{itemize}

\section{Классическая модель прецессирующего магнитного диполя}

\subsection{Внутренняя структура}

Электрон моделируется как магнитный диполь с эффективным моментом инерции $I_{\text{eff}}$, совершающий прецессию с угловой скоростью $\Omega$. Магнитный момент $\boldsymbol{\mu}$ связан с кинетическим моментом прецессии:

\begin{equation}
\boldsymbol{\mu} = \gamma \boldsymbol{L}_{\text{prec}}, \quad \boldsymbol{L}_{\text{prec}} = I_{\text{eff}} \boldsymbol{\Omega}
\end{equation}

где $\gamma$ — гиромагнитное отношение.

\subsection{Векторный потенциал электростатического поля}

Ключевым инсайтом данной модели является то, что электростатическое поле точечного заряда обладает векторным потенциалом, связанным с вращением магнитных зарядов:

\begin{equation}
\vec{A}_E = -\frac{\cot\theta}{r} \, \vec{e}_\varphi
\end{equation}

Этот векторный потенциал связан с внутренним вращением магнитных зарядов внутри структуры электрона и играет ключевую роль в стабилизации прецессии.

\section{Динамика в неоднородном магнитном поле}

\subsection{Уравнения движения}

Движение прецессирующего диполя в неоднородном поле $\boldsymbol{B}(\boldsymbol{r})$ описывается уравнениями:

\begin{equation}
\frac{d\boldsymbol{\mu}}{dt} = \gamma \boldsymbol{\mu} \times \boldsymbol{B}(\boldsymbol{r})
\end{equation}

\begin{equation}
\frac{d^2\boldsymbol{r}}{dt^2} = \frac{\mu}{m} \nabla(B \cos\theta)
\end{equation}

где $\theta$ — угол между $\boldsymbol{\mu}$ и $\boldsymbol{B}$.

\subsection{Эффективный потенциал}

Полная энергия системы включает как потенциальную энергию в магнитном поле, так и кинетическую энергию прецессии:

\begin{equation}
E_{\text{total}} = -\mu B \cos\theta + \frac{1}{2} I_{\text{eff}} \Omega^2 \sin^2\theta
\end{equation}

Первый член представляет энергию магнитного взаимодействия, второй член — кинетическую энергию прецессии.

\subsection{Стационарные состояния}

Условие $\frac{dE_{\text{total}}}{d\theta} = 0$ даёт:

\begin{equation}
\sin\theta \left( \mu B - I_{\text{eff}} \Omega^2 \cos\theta \right) = 0
\end{equation}

Это уравнение имеет решения:
\begin{itemize}
\item $\theta = 0, \pi$ (параллельное и антипараллельное состояния)
\item $\cos\theta = \frac{\mu B}{I_{\text{eff}} \Omega^2}$ (дополнительные состояния при $|\alpha| < 1$)
\end{itemize}

где $\alpha = \frac{\mu B}{I_{\text{eff}} \Omega^2}$.

\section{Связь с опытом Штерна-Герлаха}

\subsection{Стандартная конфигурация}

В традиционной постановке опыта Штерна-Герлаха градиент магнитного поля параллелен направлению поля. Сила, действующая на диполь:

\begin{equation}
F_z = \mu_z \frac{\partial B_z}{\partial z} = \mu \cos\theta \frac{\partial B_z}{\partial z}
\end{equation}

Пучок расщепляется на компоненты, соответствующие дискретным значениям $\cos\theta$.

\subsection{Перпендикулярная конфигурация}

Модифицированная конфигурация с перпендикулярными полем и градиентом ($\boldsymbol{B} \parallel z$, $\nabla\boldsymbol{B} \parallel x$) представляет фундаментальный вопрос: что определяет ось измерения?

\begin{itemize}
\item Направление поля $\boldsymbol{B}$ определяет ось квантования
\item Направление градиента $\nabla\boldsymbol{B}$ определяет ось пространственного разделения
\end{itemize}

В стандартных опытах эти оси совпадают. Однако в перпендикулярной конфигурации они различаются, что вызывает концептуальные вопросы о природе квантового измерения.

\subsection{Объяснение двух пятен}

Классическая модель объясняет два дискретных пятна через два устойчивых стационарных состояния:

\begin{itemize}
\item $\theta = 0$: Параллельное состояние (минимум энергии)
\item $\theta = \pi$: Антипараллельное состояние (стабилизировано гироскопическим эффектом прецессии)
\end{itemize}

Кинетический момент прецессии обеспечивает гироскопическую стабилизацию, которая предотвращает переворот антипараллельного состояния, в отличие от простых классических моделей, где оно было бы неустойчиво.

\section{Соответствие с квантовой механикой}

\subsection{Матрицы Паули}

Классическая модель прецессии устанавливает соответствие с формализмом матриц Паули:

\begin{center}
\begin{tabular}{|c|c|c|}
\hline
\textbf{Матрица Паули} & \textbf{Физический смысл} & \textbf{Классический аналог} \\
\hline
$\sigma_z$ & Проекция спина на ось $z$ & Угол $\theta$ прецессии \\
$\sigma_x, \sigma_y$ & Перпендикулярные компоненты & Фаза $\varphi$ прецессии \\
$\sigma^2 = 3$ & Квадрат спина & Интенсивность прецессии \\
\hline
\end{tabular}
\end{center}

\subsection{Уравнение Дирака}

Эффективный гамильтониан для прецессирующего диполя:

\begin{equation}
H_{\text{eff}} = \frac{p^2}{2m} - \boldsymbol{\mu} \cdot \boldsymbol{B} + \frac{1}{2} I_{\text{eff}} \Omega^2 \sin^2\theta + q\phi
\end{equation}

соответствует гамильтониану Дирака при квантовании. Матрицы $\alpha$ и $\beta$ находят интерпретацию:

\begin{itemize}
\item $\alpha^i$: Операторы пространственного поворота диполя
\item $\beta$: Оператор переворота диполя (север $\leftrightarrow$ юг)
\end{itemize}

\subsection{Соответствие параметров}

\begin{center}
\begin{tabular}{|c|c|c|}
\hline
\textbf{Квантовый параметр} & \textbf{Классический аналог} & \textbf{Физический смысл} \\
\hline
$\hbar/2$ & $I_{\text{eff}} \Omega$ & Квант действия \\
$\mu_B$ & $\mu$ & Магнитный момент \\
$\gamma$ & $\mu/L_{\text{prec}}$ & Гиромагнитное отношение \\
\hline
\end{tabular}
\end{center}

\section{Дополнительные стационарные состояния}

\subsection{Бифуркационная диаграмма}

Параметр $\alpha = \frac{\mu B}{I_{\text{eff}} \Omega^2}$ определяет количество стационарных состояний:

\begin{center}
\begin{tabular}{|c|c|c|}
\hline
$\alpha$ & \textbf{Стационарные состояния} & \textbf{Физическая интерпретация} \\
\hline
$\alpha > 1$ & $\theta = 0$ & Только параллельное состояние \\
$0 < \alpha < 1$ & $\theta = 0, \arccos\alpha, \pi - \arccos\alpha$ & Три устойчивых состояния \\
$\alpha = 0$ & Все $\theta$ & Изотропия (нет поля) \\
$-1 < \alpha < 0$ & $\theta = \pi, \arccos\alpha, 2\pi - \arccos\alpha$ & Три устойчивых состояния \\
$\alpha < -1$ & $\theta = \pi$ & Только антипараллельное состояние \\
\hline
\end{tabular}
\end{center}

\subsection{Экспериментальные следствия}

При $|\alpha| < 1$ на детекторе опыта Штерна-Герлаха должны появиться четыре различных пятна:

\begin{itemize}
\item $\theta = 0$: $\mu_z = +\mu$
\item $\theta = \arccos\alpha$: $\mu_z = +\mu\alpha$
\item $\theta = \pi - \arccos\alpha$: $\mu_z = -\mu\alpha$
\item $\theta = \pi$: $\mu_z = -\mu$
\end{itemize}

В традиционных опытах параметры были таковы, что $|\alpha| \approx 1$, и промежуточные состояния сливаются с основными.

\section{Анализ устойчивости}

Вторая производная энергии определяет устойчивость:

\begin{equation}
\frac{d^2E}{d\theta^2} = \mu B \cos\theta - I_{\text{eff}}\Omega^2 \cos 2\theta
\end{equation}

Для промежуточных состояний ($\cos\theta = \alpha$):

\begin{equation}
\frac{d^2E}{d\theta^2} = I_{\text{eff}}\Omega^2 (1 - \alpha^2) > 0 \quad \text{при} \quad |\alpha| < 1
\end{equation}

Таким образом, промежуточные состояния устойчивы, когда они существуют.

\section{Обсуждение}

\subsection{Преимущества модели}

\begin{itemize}
\item Предоставляет механистическое объяснение спина без сверхсветовых скоростей
\item Объясняет оба состояния (параллельное и антипараллельное) через гироскопическую стабилизацию
\item Предсказывает дополнительные состояния, проверяемые в слабопольных экспериментах
\item Устанавливает чёткое соответствие с квантовым формализмом
\item Предлагает интуитивную картину векторного потенциала и геометрической фазы
\end{itemize}

\subsection{Открытые вопросы}

\begin{itemize}
\item Определение $I_{\text{eff}}$ и $\Omega$ из первых принципов
\item Связь с релятивистскими эффектами и $Zitterbewegung$
\item Объяснение античастиц и аннигиляции
\item Квантование угла прецессии $\theta$
\end{itemize}

\section{Заключение}

Модель прецессирующего магнитного диполя предлагает классическую интерпретацию спина электрона, разрешающую фундаментальные концептуальные трудности. Рассматривая электрон как структурированный объект с внутренней прецессией, модель объясняет:

\begin{itemize}
\item Дискретные результаты опытов Штерна-Герлаха
\item Устойчивость обоих состояний (параллельного и антипараллельного)
\item Соответствие с квантово-механическим формализмом
\item Существование дополнительных стационарных состояний
\end{itemize}

Этот подход устраняет разрыв между классической интуицией и квантовым формализмом, предполагая, что спин может иметь более глубокое механическое происхождение, чем считалось ранее. Экспериментальная проверка дополнительных состояний в слабопольных конфигурациях опыта Штерна-Герлаха могла бы предоставить решающие доказательства в пользу данной модели.

\section*{Благодарности}

Автор выражает благодарность за ценные обсуждения структуры электрона и векторного потенциала электростатического поля.

\bibliographystyle{plain}
\bibliography{references}

\end{document}